\section{Reduction of the exponent range}
\label{sec:supports}

Write \(P_m=P_{m,a,d}\) for the normalized moment polynomial and
\(R_m(x)=x^mP_m(1/x)\) for its monic reciprocal. Put \(b=a+d\) and
\(s=d^2/b\). For \(x\in\mathbb R\) and even \(m\), the Gamma integral gives
\[
 R_m(x)=\frac{1}{(am)!}
       \int_0^\infty u^{am}(x+u^d)^m e^{-u}\,du>0 .
\]
This observation will also be used in the arithmetic argument.

\begin{lemma}[Root annuli and separation of the shapes]
\label{lem:gamma-annulus}
For \(a\ge1\), every root \(x\) of \(R_m\) satisfies
\[
 am<|x|^{1/d}\le (a+d)m .
\]
For any \(a\ge0\), the pair \(P_{N-1},P_N\) is coprime if
\(a\ge d(N-2)\). If \(N\ge25\), it is also coprime if
\(6a\ge (5N-11)d\).
\end{lemma}

The proof is in Appendix~\ref{sec:gamma-annulus-proof}. It uses the
root annuli of the unit-step Gamma polynomials and the fact that
these annuli multiply under composition in the binomial basis.
The strengthened inner radius matters because the factorial arguments
in the two moments grow at different rates.

\begin{lemma}[Constant plus a power]
\label{lem:constant-support}
For \(a=0\) and every \(d\ge1\), the consecutive polynomials
\(P_{N-1,0,d},P_{N,0,d}\) are coprime over \(\mathbb Q\).
\end{lemma}

\begin{proof}
In this case \(R_N'=NR_{N-1}\), by coefficient comparison. The
squarefreeness proof in Appendix~\ref{app:constant-proof} therefore
gives the result.
\end{proof}

\begin{lemma}[Endpoint normalization]
\label{lem:center-ratio}
For \(a\ge1\) define
\[
 G_m=\left(\frac{((a+d)m)!}{(am)!}\right)^{1/m},
 \qquad
 P_m(w/G_m)=1+w^m+\sum_{k=1}^{m-1}c_{m,k}w^k .
\]
For \(1\le k<m\) one has
\begin{equation}\label{eq:normalized-coefficients}
 0<c_{m,k}\le
 \binom{m}{k}\exp\!\left(-\frac{s\,k(m-k)}{2m}\right).
\end{equation}
For \(N\ge2\) the centers satisfy
\[
 G_N/G_{N-1}>(N/(N-1))^{d/2}.
\]
\end{lemma}

\begin{proof}
The consecutive differences of \(\log((am+dk)!)\) are the logarithms
of rising products of length \(d\). For \(1\le k<m\) the second
difference is
\[
 \sum_{j=1}^d\log\!\left(1+\frac{d}{am+d(k-1)+j}\right)
 \ \ge\ \frac{d^2}{(a+d)m}.
\]
Indeed, \(\log(1+v)\ge v/(1+v)\) for \(v\ge0\). Subtracting the
quadratic with this second difference and comparing with the chord
through the endpoints proves \eqref{eq:normalized-coefficients}.

For the second claim, let \(f(u)=\log\Gamma(u+1)\) and
\(H_0(x)=(f((a+d)x)-f(ax))/x\). The trigamma series gives
\(f''(u)>1/(u+1)\), by comparison with the integral. Since
\(a x\ge1\) for \(x\ge1\),
\[
 H_0'(x)=\frac{1}{x^2}\int_{ax}^{(a+d)x}u f''(u)\,du
       >\frac{d}{2x}.
\]
Integrate from \(N-1\) to \(N\).
\end{proof}

\begin{lemma}[Ratio bound]
\label{lem:ratio-range}
For \(a\ge1,d\ge1,N\ge25\), if \(s\ge\sigma_N\), where
\[
 \sigma_N=
 \begin{cases}
  2\log N+3& (25\le N<200),\\
  2\log N+2& (N\ge200),
 \end{cases}
\]
then \(P_{N-1},P_N\) are coprime.
\end{lemma}

\begin{proof}
The upper bound on the
normalized coefficients in \eqref{eq:normalized-coefficients}
is symmetric in \(k,m-k\).

Suppose \(s\ge\sigma_N\). Let \(C\) be the constant term \(3\) or \(2\)
in \(\sigma_N\), and use the following bounds for \(m=N-1,N\).
\[
\begin{array}{c|cccc}
 N&A&E&q&\eta\\ \hline
 25\le N<200&.224&1.218&.52&7/8\\
 N\ge200&.368&1.033&.24&19/20
\end{array}
\]
For \(1\le k\le4\), the right side of
\eqref{eq:normalized-coefficients} is at most
\((AE^k)^k/k!\). Here \(A\ge e^{-C/2}\) and
\(E\ge\exp((\log N+C/2)/(N-1))\).
For \(5\le k\le m/2\), put \(u=k/m\) and use
\(\binom mk\le(e/u)^k\). The logarithm of the resulting bound,
divided by \(k\), is \(1-\log u-\sigma_N(1-u)/2\).
By convexity its maximum is at an endpoint.
The tabulated \(q\) bounds both endpoint exponentials, which
are at most \(eAE^5/5\) and \(2e\,e^{-\sigma_N/4}\).
The bounds on \(E\) and on the second endpoint are checked
at the smallest \(N\) in each row, where they are largest.
In both rows
\[
 2\sum_{k=1}^{4}\frac{(AE^k)^k}{k!}+\frac{2q^5}{1-q}<\eta .
\]
Consequently \(\sum_{k=1}^{m-1}c_{m,k}<\eta\). At any root
of \(P_m(w/G_m)\) with \(|w|\le1\), the constant term gives
\(|w|^m>1-\eta\). The reciprocal polynomial gives
\(|w|^m<(1-\eta)^{-1}\) when \(|w|\ge1\).

By Lemma~\ref{lem:center-ratio} the two annuli are disjoint since
\[
 \frac d2\log\!\left(\frac N{N-1}\right)
              \frac{N(N-1)}{2N-1}
 \ \ge\ \frac{d(N-1)}{2(2N-1)}
 \ >\ -\log(1-\eta).
\]
For the last inequality, \(d>s\ge\sigma_N\) gives \(d\ge10\)
in the first row and \(d\ge13\) in the second. The minimum
of the displayed rational lower bound in each row occurs at
its smallest \(N\) and smallest \(d\). This contradicts
a common root and proves the assertion.
\end{proof}


\begin{lemma}[Support reduction]
\label{lem:support-reduction}
If \(N\ge25\) and \(P_{N-1,a,d},P_{N,a,d}\) have a common root, then
\[
 a\ge5,\qquad 2\le d<3a,\qquad
 a<d(N-2),\qquad 6a<(5N-11)d,\qquad
 s<\sigma_N,
\]
where \(s=d^2/(a+d)\).
\end{lemma}

\begin{proof}
Lemma~\ref{lem:constant-support} disposes of \(a=0\).
Lemma~\ref{lem:positive-boundary} covers \(a=1,d\ge3\) and
\(a\ge2,d\ge3a\). The remaining small positive
lower exponents and the adjacent exponents are dealt with in
Proposition~\ref{prop:additional-boundary}. These exclusions give
\(a\ge5\) and \(2\le d<3a\). Lemma~\ref{lem:gamma-annulus} gives
the two bounds on the shape. Lemma~\ref{lem:ratio-range} gives the
bound on \(s\).
\end{proof}
